\section{Background and Problem Setting}
\label{sec:background}

\subsection{Schedules dominate performance}

Since the publication of Halide~\citep{ragankelley2012halide}, it has become standard to treat the \emph{schedule} of a computation, the choices of how it is tiled, parallelized, and pipelined on the hardware, as a first-class optimization object, separate from the \emph{algorithm} that specifies what is computed. A substantial body of subsequent work has shown that these schedule choices can dominate performance: the same computation, compiled from the same source, can run several times faster or slower depending on how it is scheduled. Halide autoschedulers, TVM, AutoTVM, Ansor, TensorIR, and MetaSchedule have developed increasingly sophisticated ways to search this space, using both handcrafted heuristics and learned cost models~\citep{mullapudi2016halide,adams2019halide,chen2018tvm,chen2018learningtensor,zheng2020ansor,zheng2022tensorir,shao2022metaschedule}. The broad lesson behind these systems is that, even when the high-level computation is fixed, a substantial fraction of end-to-end performance is determined by a relatively small set of scheduling decisions.

\subsection{Triton exposes a bounded tuning surface}

Triton provides a lower-level, but still schedule-oriented, environment for writing GPU kernels~\citep{tillet2019triton,openai2021triton}. Rather than describing a high-level algorithm and asking a compiler to derive a schedule, a Triton kernel is written directly in terms of block-level operations, and the developer is expected to choose a small number of parameters that determine how the kernel is executed on the device. For matrix multiplication, these include the tile dimensions used to decompose the output matrix, the grouping used when launching thread blocks, the number of warps per block, and the number of pipeline stages used to overlap memory and compute. For layer normalization they include a block size and, again, the launch parameters.

Although this tuning surface is small compared to the full schedule language of a tensor compiler such as TVM, it is still large enough to matter. The best tile shape for a matrix multiplication depends on the input dimensions; the best number of warps depends on the register pressure of the generated code; the best number of pipeline stages depends on the memory system of the specific GPU. A developer who simply takes the default values often leaves a significant fraction of performance on the table, while a developer who benchmarks every candidate configuration quickly runs out of time.

\begin{table}[t]
\centering
\small
\begin{tabular}{p{0.21\columnwidth}p{0.69\columnwidth}}
\toprule
Item & Final choice \\
\midrule
Research question & Can cheap compile-adjacent signals and bounded profiling guide Triton schedule selection better than default or naive matched-budget baselines? \\
Primary kernel & Representative GEMM \\
Secondary kernel & LayerNorm, split into \texttt{small\_batch} and \texttt{large\_batch} regimes \\
Environment & Homogeneous RTX A6000 pool across \texttt{gpunode2}/\texttt{gpunode3} \\
Headline result & Bounded positive non-\texttt{split\_k} mainline result via \texttt{v5\_mainline\_profiled} \\
Closeout rule & No further selector-family expansion is justified in the paper backbone \\
\bottomrule
\end{tabular}
\caption{Study scope and final research framing.}
\label{tab:scope}
\end{table}

\begin{table}[t]
\centering
\small
\begin{tabular}{p{0.22\columnwidth}p{0.34\columnwidth}p{0.32\columnwidth}}
\toprule
Evaluation surface & Active tuning knobs & Evidence used by the selector \\
\midrule
Matrix multiplication (final) & Output tile dimensions (in three axes), thread-block grouping factor, warps per block, pipeline depth & Compile-adjacent signals (register count, shared-memory footprint, occupancy estimate) plus a small, well-attributed runtime profiling tier \\
Layer normalization (final) & Block size, warps per block, pipeline depth & Compile-adjacent signals plus a small layer-normalization profiling tier \\
Retired from final surface & Reduction-splitting parameter (matrix multiplication), row-grouping parameter (layer normalization) & Retained only as diagnostic evidence explaining why each parameter was removed from the main evaluation \\
\bottomrule
\end{tabular}
\caption{Tuning knobs and evidence tiers on the final evaluation surface. The knobs marked ``retired'' were evaluated earlier in the project but are no longer part of the main tuning surface; their retirement is explained in Section~\ref{sec:results}.}
\label{tab:knobs}
\end{table}


\subsection{Evidence under a budget}

The core question we study in this paper is not whether a good schedule exists, it almost always does, but whether \emph{cheap evidence} can guide us to it without exhausting the measurement budget. To make the question precise, we distinguish three kinds of evidence that can be gathered about a candidate schedule.

The first is \emph{admissibility}. Some configurations are invalid on their own terms, for example because they would require more shared memory than the device provides, or because their tile dimensions do not divide the input shape in a supported way. Rejecting these candidates is essentially free.

The second is \emph{compile-adjacent evidence}: information that can be extracted after the kernel is compiled but without actually running it on data. This includes the number of registers the kernel uses, the amount of shared memory it allocates, and an occupancy estimate that can be derived from these quantities. This kind of evidence is cheap to collect, but it describes the code rather than its actual runtime behavior.

The third is \emph{runtime evidence}: either a direct benchmark of the kernel on representative inputs, or profiling counters collected while the kernel runs. This is the most faithful kind of evidence, but also the most expensive, and it is limited by the total measurement budget.

A matched-budget tuning study asks which of these kinds of evidence are worth paying for, under a fixed total budget, and whether the resulting schedule is measurably better than the one chosen by defaults or by a naive random search spending the same budget on raw benchmark runs. That is the question we study throughout the rest of the paper.

\subsection{Workloads and their roles}

A tuning study is only as informative as the workloads on which it is evaluated. We study two Triton kernels, and we assign them deliberately different roles. Matrix multiplication is our primary workload. We evaluate it on two sets of shapes: a \emph{representative} set whose dimensions are chosen to mirror the irregular aspect ratios and masked edges that appear in real transformer workloads, and an \emph{aligned} set whose dimensions are convenient powers of two. The aligned set is retained in the paper not as the main truth source, but as a warning: as we will show in Section~\ref{sec:results}, a selector that looks strong on convenient shapes can be substantially weaker on representative ones. Layer normalization is a secondary workload, used to test whether the same tuning strategy generalizes across kernels with a different compute-to-memory profile, and its results are split by input regime so that its behavior can be interpreted rather than averaged away.

\begin{table}[t]
\centering
\small
\setlength{\tabcolsep}{4pt}
\begin{tabularx}{\columnwidth}{
  >{\raggedright\arraybackslash}p{0.21\columnwidth}
  >{\raggedright\arraybackslash}X
  >{\raggedright\arraybackslash}p{0.31\columnwidth}}
\toprule
Kernel & Workload classes & Role \\
\midrule
Representative GEMM &
\texttt{square\_compute}, \texttt{m\_dominant}, \newline
\texttt{n\_dominant}, \texttt{edge\_nondivisible} &
Primary truth source for ranking quality and the final headline \\
Aligned GEMM &
\texttt{aligned\_square}, \texttt{aligned\_rectangular} &
Context workload used to check whether regular shapes flatter the selector \\
LayerNorm &
\texttt{small\_batch}, \texttt{large\_batch} &
Secondary regime-split case for interpreting profiling value \\
\bottomrule
\end{tabularx}
\caption{Final workload matrix used by the draft.}
\label{tab:workloads}
\end{table}


\subsection{Scope and boundaries}

This paper is narrower than a full autoscheduling system. We keep the Triton implementation of each kernel fixed and study the effect of tuning only its exposed knobs under an explicit budget. Broader compiler-level synthesis and agentic systems that generate or rewrite entire CUDA kernels~\citep{dai2026cudaagent} are outside our scope. Our goal is to describe, as cleanly as possible, where cheap evidence helps, where it falls short, and what should be removed from the evaluation when a knob family fails to transfer between search spaces.
